\section{Selección del Transistor de Paso, Verificación de SOA y Cálculo de Disipador}

\subsection{Determinación de las Exigencias Eléctricas y Térmicas sobre $Q_1$}
De acuerdo con la metodología establecida en las Secciones 6-8-3 y 6-9-5h del material de cátedra \cite{gonzalez_p2, gonzalez_p3}, la selección del transistor bipolar de paso serie $Q_1$ se rige por tres parámetros fundamentales de trabajo:
\begin{enumerate}
    \item \textbf{Tensión colector-emisor máxima ($V_{CE,\max}$):} En régimen de operación normal a $V_{in} = \SI{25,0}{\volt}$, siendo $V_{i,\text{reg}} \approx \SI{9,33}{\volt}$ (\texttt{Circuito B}) a \SI{9,54}{\volt} (\texttt{Circuito A}), la tensión soportada por el transistor (\SI{15,46}{\volt} en \texttt{Circuito A}) es:
    \begin{equation}
        V_{CE,\text{op}} = V_{in} - V_{i,\text{reg}} = \SI{25,0}{\volt} - \SI{9,33}{\volt} = \SI{15,67}{\volt}
        \label{eq:vce_op}
    \end{equation}
    En condición transitoria de encendido (con el capacitor $C_1$ inicialmente descargado) o ante un eventual cortocircuito en la entrada del regulador con sobretensión de línea del $+\SI{10}{\percent}$, la tensión máxima absoluta que puede aparecer entre colector y emisor (Ec. 6-77 de \cite{gonzalez_p2}) asciende a:
    \begin{equation}
        V_{CE,\max(\text{abs})} = 1,10 \cdot V_{in} = \SI{27,5}{\volt}
        \label{eq:vce_max_abs}
    \end{equation}

    \item \textbf{Corriente de colector máxima ($I_{C,\max}$):} Despreciando la corriente de base frente a la de colector ($I_C \approx I_E = I_L + I_d$), se evalúan los tres regímenes operativos de interés:
    \begin{itemize}
        \item Régimen del Ejemplo 6-9-5h de cátedra ($I_L = \SI{300}{\milli\ampere}$): $I_{C(300)} = \SI{304,2}{\milli\ampere}$.
        \item Régimen nominal de hoja de datos del \texttt{LM7805C} ($I_{L,\text{nom}} = \SI{500}{\milli\ampere}$): $I_{C,\text{nom}} = \SI{504,2}{\milli\ampere}$.
        \item Régimen de plena carga continua del \texttt{LM7805C} ($I_{L,\max} = \SI{1,0}{\ampere}$): $I_{C,\max} = \SI{1,005}{\ampere}$ (y hasta $I_{SC} \approx \SI{1,5}{\ampere}$ durante el límite de corriente interno).
    \end{itemize}

    \item \textbf{Potencia máxima disipada en el colector ($P_{C,\max}$):} Aplicando la Ec. 6-63 de \cite{gonzalez_p2} ($P_C = V_{CE} \cdot I_C$) en $V_{in} = \SI{25,0}{\volt}$ (\SI{7,79}{\watt} a \SI{500}{\milli\ampere} en \texttt{Circuito A}):
    \begin{align}
        P_{C(300\text{mA})} &= \SI{15,67}{\volt} \times \SI{0,3042}{\ampere} = \SI{4,77}{\watt} \label{eq:pc_300}\\
        P_{C,\text{nom}(500\text{mA})} &= \SI{15,67}{\volt} \times \SI{0,5042}{\ampere} = \SI{7,90}{\watt} \label{eq:pc_500}\\
        P_{C,\max(1,0\text{A})} &= \SI{15,67}{\volt} \times \SI{1,005}{\ampere} = \SI{15,75}{\watt} \label{eq:pc_1000}
    \end{align}
\end{enumerate}

\subsection{Selección Comparativa del Transistor}
En la Tabla~\ref{tab:comparativa_transistores} se contrastan los dos candidatos estándar empleados en los apuntes de la cátedra: el transistor de media potencia \textbf{\texttt{BD135-16}} (incluido en el Apéndice 2 de \cite{gonzalez_p2}) y el transistor de potencia \textbf{\texttt{TIP31C}} (familia \texttt{TIP31A/B/C}, empleado en la Sección 6-9-5h de \cite{gonzalez_p3}).

\begin{table}[htbp]
\centering
\caption{Comparación de especificaciones límite y térmicas entre los transistores \texttt{BD135-16} y \texttt{TIP31C} \cite{gonzalez_p2, onsemi_tip31}.}
\label{tab:comparativa_transistores}
\resizebox{\columnwidth}{!}{%
\begin{tabular}{llccc}
\toprule
\textbf{Parámetro} & \textbf{Símbolo} & \textbf{Requerido ($I_L = \SI{0,5}{\ampere} / \SI{1,0}{\ampere}$)} & \textbf{\texttt{BD135-16} (TO-126)} & \textbf{\texttt{TIP31C} (TO-220)} \\
\midrule
Tensión colector-emisor & $V_{CEO(\text{sus})}$ & $> \SI{27,5}{\volt}$ & \SI{45}{\volt} & \textbf{\SI{100}{\volt}} (\SI{60}{\volt} en \texttt{TIP31A}) \\
Tensión colector-base & $V_{CBO}$ & $> \SI{27,5}{\volt}$ & \SI{45}{\volt} & \textbf{\SI{100}{\volt}} \\
Corriente de colector (DC) & $I_{C(\text{DC})}$ & \SI{0,504}{\ampere} / \SI{1,005}{\ampere} & \SI{1,5}{\ampere} & \textbf{\SI{3,0}{\ampere}} \\
Corriente de colector pico & $I_{CM}$ & $> \SI{1,5}{\ampere}$ & \SI{3,0}{\ampere} & \textbf{\SI{5,0}{\ampere}} \\
Potencia máx. ($T_C = \SI{25}{\celsius}$) & $P_{C(T_C)}$ & \SI{7,90}{\watt} / \SI{15,75}{\watt} & \SI{12,5}{\watt} (\emph{insuficiente a \SI{1}{\ampere}}) & \textbf{\SI{40,0}{\watt}} \\
Potencia máx. ($T_A = \SI{25}{\celsius}$) & $P_{C(T_A)}$ & --- & \SI{1,25}{\watt} & \textbf{\SI{2,0}{\watt}} \\
Ganancia de corriente DC & $h_{FE}$ & $\ge 50$ & $100\text{--}250$ & $\mathbf{50\text{--}100}$ ($\approx 90$ típ.) \\
Temp. de juntura máxima & $T_{J,\max}$ & $> \SI{125}{\celsius}$ & \SI{150}{\celsius} & \textbf{\SI{150}{\celsius}} \\
Res. térmica juntura-carcasa & $R_{th(j-mb)}$ & Baja para disipador & \SI{10,0}{\celsius\per\watt} & \textbf{\SI{3,125}{\celsius\per\watt}} \\
Res. térmica juntura-ambiente & $R_{th(j-a)}$ & --- & \SI{100,0}{\celsius\per\watt} & \textbf{\SI{62,5}{\celsius\per\watt}} \\
\bottomrule
\end{tabular}%
}
\end{table}

\textbf{Justificación de la elección:} Si bien el \texttt{BD135-16} podría utilizarse para corrientes reducidas ($I_L \le \SI{300}{\milli\ampere}$), para $I_L = \SI{500}{\milli\ampere}$ su elevada resistencia térmica interna ($R_{th(j-mb)} = \SI{10,0}{\celsius\per\watt}$) produce por sí sola una caída térmica juntura-carcasa de $\Delta T_{j-mb} = \SI{79,0}{\celsius}$, y a plena capacidad del regulador ($I_L = \SI{1,0}{\ampere}$, $P_C = \SI{15,75}{\watt}$) supera la potencia máxima admisible de su encapsulado TO-126 (\SI{12,5}{\watt}). En contrapartida, el transistor \textbf{\texttt{TIP31C}} (en encapsulado TO-220) ofrece $P_{C,\max} = \SI{40,0}{\watt}$ (más de $5\times$ la potencia nominal de \SI{7,90}{\watt} y $2,5\times$ la de plena carga de \SI{15,75}{\watt}), $I_{C,\max} = \SI{3,0}{\ampere}$, $V_{CEO} = \SI{100}{\volt}$ y una resistencia térmica juntura-base de montaje más de tres veces menor ($R_{th(j-mb)} = \SI{3,125}{\celsius\per\watt}$). Por lo tanto, \textbf{se selecciona el transistor \texttt{TIP31C}}.

\subsection{Verificación del Punto de Operación según el Área de Operación Segura (SOA)}
Conforme a la Sección 6-8-3c del material teórico \cite{gonzalez_p2}, el diagrama del Área de Operación Segura (\emph{Safe Operating Area}, SOA) representa en ejes logarítmicos la corriente de colector $I_C$ en función de la tensión colector-emisor $V_{CE}$, delimitando la región de trabajo exenta de fallo mediante cuatro fronteras físicas:
\begin{enumerate}
    \item \textbf{Frontera (a) -- Límite de corriente máxima en continua ($I_C = \SI{3,0}{\ampere}$):} Recta horizontal impuesta por la sección transversal de los hilos de interconexión interna (\emph{bonding wires}) y la densidad máxima de portadores.
    \item \textbf{Frontera (b) -- Límite de disipación térmica máxima ($P_{C,\max} = \SI{40,0}{\watt}$ a $T_C = \SI{25}{\celsius}$):} Hipérbola de potencia constante $I_C = P_{C,\max} / V_{CE}$, que en representación doble logarítmica constituye una recta de pendiente $-1$ (\SI{-20}{\deci\bel}/década) entre $V_{CE} = \SI{13,33}{\volt}$ y $V_{CE} = \SI{40,0}{\volt}$, gobernada por la resistencia térmica $R_{th(j-mb)}$.
    \item \textbf{Frontera (c) -- Límite de segunda ruptura (\emph{Second Breakdown}, $V_{CE} > \SI{40}{\volt}$):} Recta de pendiente más abrupta ($\approx -1,65$) ocasionada por la focalización térmica de la corriente en puntos calientes microscópicos de la juntura base-colector bajo campos eléctricos intensos, que llevaría a la fusión local del silicio si no se reduce la corriente admisible.
    \item \textbf{Frontera (d) -- Límite de ruptura por avalancha ($V_{CEO(\text{sus})} = \SI{100}{\volt}$):} Recta vertical que establece la tensión máxima colector-emisor con base abierta.
\end{enumerate}

En la Figura~\ref{fig:soa_plot} se grafica la curva SOA del \texttt{TIP31C} tanto a $T_C = \SI{25}{\celsius}$ como aplicando la reducción térmica (\emph{derating}) para una temperatura de carcasa realista en servicio de $T_C = \SI{85}{\celsius}$ ($P_{C(\SI{85}{\celsius})} = \SI{20,8}{\watt}$), superponiendo además el límite del \texttt{BD135-16} y la trayectoria del punto de operación $Q(V_{CE}, I_C)$ para $I_L = \SI{100}{\milli\ampere}$, \SI{300}{\milli\ampere}, \SI{500}{\milli\ampere} y \SI{1,0}{\ampere}. Se comprueba visual y analíticamente que todos los puntos de operación a $V_{in} = \SI{25}{\volt}$ ($V_{CE} \approx \SI{15,5}{\volt}\text{--}\SI{15,7}{\volt}$) y a $V_{in} = \SI{10}{\volt}$ ($V_{CE} \approx \SI{1,94}{\volt}$) se ubican con amplio margen de seguridad dentro de la región SOA del \texttt{TIP31C}, incluso bajo régimen térmico de carcasa a \SI{85}{\celsius}.

\begin{figure}[htbp]
\centering
\begin{tikzpicture}
\begin{loglogaxis}[
    width=\columnwidth,
    height=7.0cm,
    xlabel={Tensión Colector-Emisor $V_{CE}$ [\si{\volt}]},
    ylabel={Corriente de Colector $I_C$ [\si{\ampere}]},
    xmin=1, xmax=180,
    ymin=0.02, ymax=5.5,
    grid=both,
    grid style={line width=.1pt, draw=gray!25},
    major grid style={line width=.2pt, draw=gray!50},
    legend pos=south west,
    legend style={font=\tiny, fill=white, fill opacity=0.92, draw=gray!60, inner sep=1.5pt, row sep=0pt},
    tick label style={font=\footnotesize},
    label style={font=\footnotesize\bfseries},
    title style={font=\footnotesize\bfseries, align=center},
    title={Verificación del Punto de Operación en el Diagrama SOA (CC)}
]
    % Curva SOA TIP31C Tc = 25 C
    \addplot[name path=soa25, color=blue!80!black, ultra thick] table[x=Vce, y=IC_TIP31C_25C] {sim/soa_data.txt};
    \addlegendentry{SOA \texttt{TIP31C} ($T_C = \SI{25}{\celsius}$, \SI{40}{\watt})}

    % Curva SOA TIP31C Tc = 85 C
    \addplot[name path=soa85, color=teal!80!black, thick, dashed] table[x=Vce, y=IC_TIP31C_85C] {sim/soa_data.txt};
    \addlegendentry{SOA \texttt{TIP31C} ($T_C = \SI{85}{\celsius}$, \SI{20,8}{\watt})}

    % Curva SOA BD135 Tc = 25 C
    \addplot[color=purple!80!black, thick, dashdotted] table[x=Vce, y=IC_BD135_25C] {sim/soa_data.txt};
    \addlegendentry{SOA \texttt{BD135-16} ($T_C = \SI{25}{\celsius}$, \SI{12,5}{\watt})}

    % Puntos de operación a Vin = 25 V (Vce = 15.5 V)
    \addplot[only marks, mark=*, mark size=2.5pt, color=red!85!black] coordinates {
        (15.46, 0.104)
        (15.50, 0.304)
        (15.55, 0.504)
        (15.67, 1.005)
    };
    \addlegendentry{Puntos $Q$ ($V_{in}=\SI{25}{\volt}$)}

    % Puntos de operación a Vin = 10 V (Vce = 1.94 V)
    \addplot[only marks, mark=square*, mark size=2.3pt, color=orange!90!black] coordinates {
        (1.94, 0.104)
        (1.94, 0.304)
        (1.94, 0.504)
    };
    \addlegendentry{Puntos $Q$ ($V_{in}=\SI{10}{\volt}$)}

    % Línea de trayectoria de carga en Vin = 25 V
    \addplot[color=red!75!black, thick, densely dotted] coordinates {(15.46, 0.104) (15.67, 1.005)};

    % Etiquetas de los puntos y fronteras
    \node[anchor=west, font=\tiny, text=red!80!black] at (axis cs:17.0, 0.104) {$I_L=\SI{0,1}{\ampere}$ (\SI{1,6}{\watt})};
    \node[anchor=west, font=\tiny, text=red!80!black] at (axis cs:17.0, 0.304) {$I_L=\SI{0,3}{\ampere}$ (\SI{4,8}{\watt})};
    \node[anchor=west, font=\tiny\bfseries, text=red!80!black] at (axis cs:17.0, 0.52) {$Q_{\text{nom}}: \SI{0,5}{\ampere}$ (\SI{7,9}{\watt})};
    \node[anchor=east, font=\tiny\bfseries, text=red!80!black] at (axis cs:14.2, 1.05) {$Q_{\max}: \SI{1,0}{\ampere}$ (\SI{15,8}{\watt})};

    \node[anchor=south, font=\tiny\bfseries, text=blue!80!black] at (axis cs:4.0, 3.05) {(a) $I_{C,\max}$};
    \node[anchor=south west, font=\tiny\bfseries, text=blue!80!black] at (axis cs:19.0, 2.0) {(b) $P_{C,\max}$};
    \node[anchor=south west, font=\tiny\bfseries, text=blue!80!black] at (axis cs:45.0, 0.90) {(c) 2da rupt.};
\end{loglogaxis}
\end{tikzpicture}
\caption{Diagrama de Área de Operación Segura (SOA) en corriente continua para el transistor \texttt{TIP31C} y ubicación de los puntos de trabajo $Q$ en $V_{in} = \SI{25}{\volt}$ y $V_{in} = \SI{10}{\volt}$.}
\label{fig:soa_plot}
\end{figure}

\subsection{Cálculo Térmico y Dimensionamiento de Disipadores}

\subsubsection{Modelo Electrotérmico (Ley de Ohm Térmica)}
Siguiendo la Sección 6-8-3b y la Tabla 6-2 del apunte de cátedra \cite{gonzalez_p2}, la transferencia de calor en régimen estacionario desde la juntura semiconductora a temperatura $T_J$ hacia el aire ambiente a temperatura $T_a$ se modela mediante el circuito serie de resistencias térmicas de la Figura~\ref{fig:circuito_termico} (Ecs. 6-79 a 6-81 de \cite{gonzalez_p2}):
\begin{equation}
    R_{th(j-a)} = \frac{T_J - T_a}{P_C} = R_{th(j-mb)} + R_{th(mb-h)} + R_{th(h-a)}
    \label{eq:ley_ohm_termica}
\end{equation}
donde las resistencias térmicas se expresan en $[\si{\celsius\per\watt}]$, siendo $R_{th(j-mb)}$ la resistencia térmica entre la juntura y la base metálica de montaje (carcasa), $R_{th(mb-h)}$ la resistencia térmica de la interfaz aislante (separador de mica más grasa siliconada/grafitada, $R_{th(mb-h)} \approx \SI{1,3}{\celsius\per\watt}$ según el Ejemplo 9 de \cite{gonzalez_p2}), y $R_{th(h-a)}$ la resistencia térmica entre el disipador (\emph{heatsink}) y el ambiente.

\begin{figure}[htbp]
\centering
\resizebox{\columnwidth}{!}{%
\begin{circuitikz}[scale=0.92, transform shape]
    \draw (0,0) node[ground]{} to[I, i=$P_C\ {[\si{\watt}]}$] (0,2.2)
          to[short, -*] (0.8,2.2) node[above=3pt]{\textbf{$T_J$}}
          to[R, l=$R_{th(j-mb)}$, a={\scriptsize Juntura--Carcasa}, -*] (3.3,2.2) node[above=3pt]{\textbf{$T_C$}}
          to[R, l=$R_{th(mb-h)}$, a={\scriptsize Mica (\SI{1,3}{\celsius\per\watt})}, -*] (6.1,2.2) node[above=3pt]{\textbf{$T_H$}}
          to[R, l=$R_{th(h-a)}$, a={\scriptsize Disipador--Amb.}, -*] (8.6,2.2) node[above=3pt]{\textbf{$T_a$}}
          to[V, v={$T_a = \SI{40}{\celsius}$}] (8.6,0) node[ground]{};
\end{circuitikz}%
}
\caption{Circuito térmico equivalente en estado estacionario basado en la analogía con la Ley de Ohm (Figura 6-44 de \cite{gonzalez_p2}).}
\label{fig:circuito_termico}
\end{figure}

\subsubsection{Cálculo del Disipador para el Transistor de Paso $Q_1$ (\texttt{TIP31C})}
\begin{enumerate}
    \item \textbf{Verificación de operación sin disipador:} Para el \texttt{TIP31C} sin disipador ($R_{th(j-a)} = \SI{62,5}{\celsius\per\watt}$), considerando una temperatura ambiente máxima de diseño $T_a = \SI{40}{\celsius}$ (Ejemplo 9 de \cite{gonzalez_p2}) y la disipación nominal $P_{C,\text{nom}} = \SI{7,90}{\watt}$ ($I_L = \SI{500}{\milli\ampere}$), la temperatura que alcanzaría la juntura según la ecuación~\eqref{eq:ley_ohm_termica} (superando ampliamente $T_{J,\max} = \SI{150}{\celsius}$) sería:
    \begin{equation}
        T_{J(\text{sin})} = \SI{62,5}{\celsius\per\watt} \times \SI{7,90}{\watt} + \SI{40}{\celsius} = \SI{533,8}{\celsius}
        \label{eq:tj_sin_disip_q1}
    \end{equation}
    Aun a $I_L = \SI{300}{\milli\ampere}$ ($P_C = \SI{4,77}{\watt}$), resultaría $T_J = \SI{338,1}{\celsius} > \SI{150}{\celsius}$. Se concluye que \textbf{el montaje de $Q_1$ sobre un disipador de calor es estrictamente obligatorio}.

    \item \textbf{Cálculo de $R_{th(h-a)}$ requerido:} Despejando $R_{th(h-a)}$ de la ecuación~\eqref{eq:ley_ohm_termica} y adoptando por confiabilidad una temperatura máxima de juntura de diseño $T_{J,\text{seg}} = \SI{105}{\celsius}$ para el régimen nominal de \SI{500}{\milli\ampere} (\SI{100}{\celsius} para \SI{300}{\milli\ampere} y \SI{125}{\celsius} para el régimen máximo de \SI{1,0}{\ampere}), con $R_{th(j-mb)} = \SI{3,125}{\celsius\per\watt}$ y $R_{th(mb-h)} = \SI{1,30}{\celsius\per\watt}$:
    \begin{equation}
        R_{th(h-a)} \le \frac{T_{J,\text{seg}} - T_a}{P_C} - \left( R_{th(j-mb)} + R_{th(mb-h)} \right)
        \label{eq:rth_ha_despeje}
    \end{equation}
    Evaluando la expresión~\eqref{eq:rth_ha_despeje} para los tres niveles de corriente de carga:
    \begin{itemize}
        \item \textbf{Para $I_L = \SI{300}{\milli\ampere}$ ($P_C = \SI{4,77}{\watt}$, $T_{J,\text{seg}} = \SI{100}{\celsius}$):}
        \begin{equation}
            R_{th(h-a)(300)} \le \frac{\SI{100}{\celsius} - \SI{40}{\celsius}}{\SI{4,77}{\watt}} - \SI{4,425}{\celsius\per\watt} = \SI{8,15}{\celsius\per\watt}
        \end{equation}
        \item \textbf{Para el régimen nominal $I_{L,\text{nom}} = \SI{500}{\milli\ampere}$ ($P_{C,\text{nom}} = \SI{7,90}{\watt}$, $T_{J,\text{seg}} = \SI{105}{\celsius}$):}
        \begin{equation}
            R_{th(h-a)(\text{nom})} \le \frac{\SI{105}{\celsius} - \SI{40}{\celsius}}{\SI{7,90}{\watt}} - \SI{4,425}{\celsius\per\watt} = \mathbf{\SI{3,80}{\celsius\per\watt}}
            \label{eq:rth_ha_500}
        \end{equation}
        \item \textbf{Para plena carga continua $I_{L,\max} = \SI{1,0}{\ampere}$ ($P_{C,\max} = \SI{15,75}{\watt}$, $T_{J,\text{seg}} = \SI{125}{\celsius}$):}
        \begin{equation}
            R_{th(h-a)(\max)} \le \frac{\SI{125}{\celsius} - \SI{40}{\celsius}}{\SI{15,75}{\watt}} - \SI{3,625}{\celsius\per\watt} = \mathbf{\SI{1,77}{\celsius\per\watt}}
            \label{eq:rth_ha_1000}
        \end{equation}
    \end{itemize}

    \item \textbf{Selección comercial del disipador para $Q_1$:} Del catálogo nacional \emph{International Aluel} (\url{disipadores.com}, referenciado en la pág. 26 de \cite{gonzalez_p2}) se selecciona:
    \begin{itemize}
        \item Para operación hasta $I_{L,\text{nom}} = \SI{500}{\milli\ampere}$: \textbf{Disipador aletado de aluminio extruido anodizado negro Aluel Art. ZD-1 / P-25} ($L = \SI{50}{\milli\meter}$, resistencia térmica garantizada $R_{th(h-a)} = \SI{3,2}{\celsius\per\watt} < \SI{3,80}{\celsius\per\watt}$), con el cual la temperatura real de juntura resulta $T_J = \SI{40}{\celsius} + \SI{7,90}{\watt} \times (\SI{3,125}{} + \SI{1,30}{} + \SI{3,20}{})\,\si{\celsius\per\watt} = \SI{100,2}{\celsius} < \SI{105}{\celsius}$.
        \item Para operación continua hasta $I_{L,\max} = \SI{1,0}{\ampere}$: \textbf{Disipador multialeta Aluel Art. ZD-4} ($L = \SI{75}{\milli\meter}$, $R_{th(h-a)} = \SI{1,6}{\celsius\per\watt} < \SI{1,77}{\celsius\per\watt}$), manteniendo $T_J = \SI{122,3}{\celsius} < \SI{125}{\celsius}$.
    \end{itemize}
\end{enumerate}

\subsubsection{Verificación Térmica y Disipador para el Regulador Integrado \texttt{LM7805C}}
Gracias a la inclusión de la etapa prereguladora ($V_{i,\text{reg}} = \SI{9,54}{\volt}$ en \texttt{Circuito A} o \SI{9,33}{\volt} en \texttt{Circuito B}), la potencia disipada por el regulador monolítico \texttt{LM7805C} a la corriente nominal de $I_{L,\text{nom}} = \SI{500}{\milli\ampere}$ se reduce drásticamente de \SI{10,11}{\watt} (conexión directa a \SI{25}{\volt}) a:
\begin{equation}
    P_{D(\text{7805})} = \left(V_{i,\text{reg}} - V_O\right) I_L + V_{i,\text{reg}} I_d = \SI{2,31}{\watt}
    \label{eq:pd_7805_nom}
\end{equation}
(y $P_{D(\text{7805})} = \SI{4,58}{\watt}$ para $I_{L,\max} = \SI{1,0}{\ampere}$). En aire libre sin plano de cobre extendido, el encapsulado TO-220 presenta $\theta_{JA} \approx \SI{50}{\celsius\per\watt}$, lo que produciría $T_J = \SI{40}{\celsius} + \SI{2,31}{\watt} \times \SI{50}{\celsius\per\watt} = \SI{155,5}{\celsius} > T_{J,\max(\text{op})} = \SI{125}{\celsius}$. Aplicando nuevamente la ecuación~\eqref{eq:rth_ha_despeje} con $\theta_{JC} = \SI{3,0}{\celsius\per\watt}$, $R_{th(mb-h)} = \SI{1,3}{\celsius\per\watt}$ y $T_{J,\text{seg}} = \SI{100}{\celsius}$ (para $I_{L,\max} = \SI{1,0}{\ampere}$, $P_D = \SI{4,58}{\watt}$; siendo $R_{th(h-a)} \le \SI{21,67}{\celsius\per\watt}$ a \SI{500}{\milli\ampere}):
\begin{equation}
    R_{th(h-a)(\text{7805})} \le \frac{\SI{100}{\celsius} - \SI{40}{\celsius}}{\SI{4,58}{\watt}} - \SI{4,30}{\celsius\per\watt} = \SI{8,80}{\celsius\per\watt}
    \label{eq:rth_ha_7805}
\end{equation}
Se adopta el disipador compacto en perfil ``U'' para TO-220 \textbf{International Aluel Art. 2125D} (dimensiones $\SI{29}{\milli\meter} \times \SI{29}{\milli\meter} \times \SI{29}{\milli\meter}$, altura \SI{40}{\milli\meter}, espesor \SI{1,5}{\milli\meter}, $R_{th(h-a)} = \SI{10,0}{\celsius\per\watt}$ ---exactamente el modelo del Ejemplo 9 de cátedra \cite{gonzalez_p2}--- para el caso nominal de \SI{500}{\milli\ampere}, con el cual $T_{J(\text{7805})} = \SI{73,0}{\celsius}$), o bien el montaje conjunto de $Q_1$ y el \texttt{LM7805C} sobre el perfil extruido \textbf{Aluel ZD-4}. En la Figura~\ref{fig:derating_termico} se representan las curvas de reducción de potencia (\emph{Power Derating}, Fig. 6-42 de \cite{gonzalez_p2}) de los dispositivos considerados.

\begin{figure}[htbp]
\centering
\begin{tikzpicture}
\begin{axis}[
    width=\columnwidth,
    height=5.6cm,
    xlabel={Temperatura de Carcasa $T_C$ [\si{\celsius}]},
    ylabel={Potencia Admisible $P_D$ [\si{\watt}]},
    xmin=0, xmax=160,
    ymin=0, ymax=45,
    grid=both,
    grid style={line width=.1pt, draw=gray!25},
    major grid style={line width=.2pt, draw=gray!50},
    legend pos=north east,
    legend style={font=\scriptsize, fill=white, fill opacity=0.9, draw=gray!60},
    tick label style={font=\footnotesize},
    label style={font=\footnotesize\bfseries},
    title style={font=\footnotesize\bfseries},
    title={Curvas de Reducción Térmica de Potencia (\emph{Power Derating})}
]
    \addplot[color=blue!80!black, ultra thick] table[x=Tc_C, y=PD_TIP31C_W] {sim/termico_data.txt};
    \addlegendentry{\texttt{TIP31C} ($R_{th(j-c)} = \SI{3,125}{\celsius\per\watt}$)}

    \addplot[color=teal!80!black, thick, dashed] table[x=Tc_C, y=PD_LM7805_W] {sim/termico_data.txt};
    \addlegendentry{\texttt{LM7805C} TO-220}

    \addplot[color=purple!80!black, thick, dashdotted] table[x=Tc_C, y=PD_BD135_W] {sim/termico_data.txt};
    \addlegendentry{\texttt{BD135-16} ($R_{th(j-c)} = \SI{10,0}{\celsius\per\watt}$)}

    % Puntos térmicos de operación con disipador seleccionado
    \addplot[only marks, mark=*, mark size=2.5pt, color=red!85!black] coordinates {
        (75.5, 7.90)
        (66.3, 2.31)
    };
    \node[anchor=south, font=\tiny\bfseries, text=red!80!black] at (axis cs:72, 8.8) {$Q_1$ (\SI{0,5}{\ampere}): $\SI{75,5}{\celsius}, \SI{7,9}{\watt}$};
    \node[anchor=west, font=\tiny, text=red!80!black] at (axis cs:68, 2.5) {\texttt{7805} (\SI{0,5}{\ampere}): $\SI{66,3}{\celsius}, \SI{2,3}{\watt}$};
\end{axis}
\end{tikzpicture}
\caption{Curvas de disipación máxima en función de la temperatura de carcasa $T_C$ y ubicación de los puntos térmicos de operación con los disipadores calculados.}
\label{fig:derating_termico}
\end{figure}
